
%(BEGIN_QUESTION)
% Copyright 2012, Tony R. Kuphaldt, released under the Creative Commons Attribution License (v 1.0)
% This means you may do almost anything with this work of mine, so long as you give me proper credit

Les og lag en disposisjon av ``ANSI/IEEE Function Number Codes''''-delen i kapitlet ``Electric Power Measurement og Control'' i læreboken {\it Lessons In Industrial Instrumentation}. Noter sidetall der viktige illustrasjoner, fotografier, ligninger, tabeller og andre relevante detaljer finnes. Forbered deg på en faglig diskusjon med lærer og medstudenter om begrepene og eksemplene i lesingen.

\underbar{file i01250}
%(END_QUESTION)




%(BEGIN_ANSWER)

Her er en liste over noen av de vanligste ANSI/IEEE-funksjonskodene:

\begin{itemize}
\item{} {\tt 50} = Momentan overstrømsvern
\item{} {\tt 51} = Tidsforsinket overstrømsvern
\item{} {\tt 52} = AC-effektbryter
\item{} {\tt 67} = Retningsbestemt overstrømsvern
\item{} {\tt 79} = Automatisk gjeninnkoblingsvern
\item{} {\tt 86} = Hjelpe-/lockout-funksjon
\item{} {\tt 81} = Frekvensvern
\item{} {\tt 87} = Differensialvern
\item{} {\tt 89} = Linjefraskiller
\end{itemize}

%(END_ANSWER)





%(BEGIN_NOTES)

ANSI/IEEE-definerte tallkoder refererer til spesifikke funksjoner i elektriske kraftsystemer. Hvert tall representerer en {\it funksjon}, ikke en bestemt {\it enhet}, slik at ett vernerelé kan implementere flere funksjoner (f.eks. 50/51 overstrøm).

\vskip 10pt

Her er en liste over noen av de vanligste ANSI/IEEE-funksjonskodene, der kursiv tekst viser abstrakte funksjoner heller enn konkrete enheter:

\begin{itemize}
\item{} {\tt 50} = {\it Momentan overstrømsvern}
\item{} {\tt 51} = {\it Tidsforsinket overstrømsvern}
\item{} {\tt 52} = AC-effektbryter
\item{} {\tt 67} = {\it Retningsbestemt overstrømsvern}
\item{} {\tt 79} = {\it Automatisk gjeninnkoblingsvern}
\item{} {\tt 86} = {\it Hjelpe-/lockout-funksjon}
\item{} {\tt 81} = {\it Frekvensvern}
\item{} {\tt 87} = {\it Differensialvern}
\item{} {\tt 89} = Linjefraskiller
\end{itemize}

\vskip 10pt

Enlinjediagram for vernerelé bruker heltrukne linjer for kontinuerlige (analoge) signalveier og stiplede linjer for diskrete (på/av) signaler eller kommunikasjon mellom reléer.

\vskip 10pt

Trip-kretsskjema viser individuelle ledere, i motsetning til enlinjeskjema som kun viser effektveier. Disse skjemaene vises ofte i par: ett for effektledere og instrumenttransformatorforbindelser til reléet, og ett for utløserkrets med relékontakter mot bryterens trip/close-spoler. En betegnelse som 51-2 starter med ANSI-kode (f.eks. 51 = tidsforsinket overstrøm) for funksjonen og avsluttes med nummer/bokstav for spesifikk reléenhet (f.eks. -2 = relé som overvåker fase 2). Bokstavsuffiks gir ytterligere detaljering (f.eks. 51-2/SI = ``seal-in''-kontakt for fase-2-reléet).











\vskip 20pt \vbox{\hrule \hbox{\strut \vrule{} {\bf Forslag til sokratisk diskusjon} \vrule} \hrule}

\begin{itemize}
\item{} Beskriv hvordan vernerelé-symbolikk i enlinjeskjema ligner ISA-symbolikk i P\&ID.
\item{} Identifiser ANSI/IEEE-funksjonskoden for {\it instantaneous overcurrent}, og gi et eksempel på en effektfeil som kan utløse denne tilstanden.
\item{} Identifiser ANSI/IEEE-funksjonskoden for {\it time overcurrent}, og gi et eksempel på en effektfeil som kan utløse denne tilstanden.
\item{} Identifiser ANSI/IEEE-funksjonskoden for {\it undervoltage}, og gi et eksempel på en effektfeil som kan utløse denne tilstanden.
\item{} Identifiser ANSI/IEEE-funksjonskoden for {\it overvoltage}, og gi et eksempel på en effektfeil som kan utløse denne tilstanden.
\item{} Identifiser ANSI/IEEE-funksjonskoden for {\it reclose}, og gi et eksempel på en effektfeil som kan utløse denne tilstanden.
\item{} Identifiser ANSI/IEEE-funksjonskoden for {\it lockout/auxiliary}, og gi et eksempel på en effektfeil som kan utløse denne tilstanden.
\item{} Identifiser ANSI/IEEE-funksjonskoden for {\it differential}, og gi et eksempel på en effektfeil som kan utløse denne tilstanden.
\item{} Sammenlign {\it effekt}-kretsskjema og {\it utløser}-kretsskjema.
\item{} Forklar hvordan seal-in-kontakten i et 51-relé virker, og hvordan latch-funksjonen resettes når effektbryteren står åpen (utløst).
\item{} Studer effekt/utløser-skjemaene i denne delen av læreboken og identifiser komponentsymboler. Hvor like eller ulike er de i forhold til vanlige elektriske skjema eller relé-ladderdiagram?
\item{} Studer effekt/utløser-skjemaene og identifiser hvordan duplisering reduseres for enklere lesbarhet.
\end{itemize}

%INDEX% Leseoppgave: Lessons In Industrial Instrumentation, ANSI/IEEE-funksjonskoder

%(END_NOTES)
